\hypertarget{appendix-89985e6a2685}{}
\section*{Binary long division}
\chapter{Binary long division}\label{appendix-89985e6a2685-binary-long-division}

\textbf{Curriculum depth:} Introduced · \textbf{First seen:} C002 · \textbf{Latest development:} C002

\section{Definition and mechanism}\label{appendix-89985e6a2685-definition-and-mechanism}

Shift dividend bits into a remainder, compare/subtract the divisor, and emit quotient bits while maintaining \(D=Qd+R\).

At each dividend bit, form remainder=2×remainder+bit. If it reaches the divisor, subtract and emit quotient bit 1; otherwise emit 0. For 13/3, quotient 4 and remainder 1 satisfy 13=4×3+1. A wider temporary may be needed for shifted remainder.

\section{Prerequisites}\label{appendix-89985e6a2685-prerequisites}

\begin{itemize}
\tightlist
\item
  \hyperref[appendix-73425b25118e]{Shift operations} --- Logical right shift fills high positions with zero; arithmetic right shift replicates the sign bit.
\item
  \hyperref[appendix-1c7df340573e]{Subtraction through addition} --- At N-bit residue width, A−B equals A+NOTₙ(B)+1.
\end{itemize}

\section{Used by}\label{appendix-89985e6a2685-used-by}

Connect later chunks here when they use this concept.

\section{Recall and next depth}\label{appendix-89985e6a2685-recall-and-next-depth}

Explain binary long division without the label, work the example above, then state the assumption whose removal would change the result.

\textbf{Pending:} restoring/non-restoring division, signed rounding, reciprocal methods, and hardware latency.

\section{Source and navigation}\label{appendix-89985e6a2685-source-and-navigation}

\href{Continuous\%20Notes/Computing\%20Foundations\%20-\%20From\%20Bits\%20to\%20C++\%20and\%20Systems\%20APIs\#Arithmetic\%20becomes\%20a\%20policy}{C002 teaching chapter} · \hyperref[appendix-fd090d982a15]{Part definitions and mechanisms} · \href{../02-Foundations-II/main.pdf\#nameddest=appendix-4ad47bdd269a}{Dependency Map (02-Foundations-II, appendix-4ad47bdd269a)}
